% --------------------------------------------------------------------
% V15 controlling source update: balanced-seven polarization frontier
% --------------------------------------------------------------------
\section*{V15 controlling update: balanced-seven polarization and signed endpoint distribution}
\addcontentsline{toc}{section}{V15 controlling update: balanced-seven polarization and signed endpoint distribution}

\begin{status}[V15 two-front firewall]
The V14 Vaughan/M\"obius prime-modulus branch is not withdrawn:
\[
 \texttt{AFFINE287-PRIME-MODULUS-MU-TWOOUTER45}
\]
remains an open source-exact analytic theorem on the M/Vaughan side.
The present V15 update concerns the balanced-seven V-branch.  It reduces that branch to a
different exact endpoint-distribution residual and a comparison/source pin.  No implication
between the two open branches is asserted here.
\end{status}

\subsection*{Balanced-seven literal source}

Set
\[
 Q=X^{3/5},\qquad Y=X^{1/7},\qquad s\in\{-1,+1\},
\]
and let
\[
 a_s\equiv -s\overline 2\pmod q.
\]
The balanced-seven source cell is
\[
 \mathcal B_{7,s}(X)
 =
 \sum_{q\sim Q}\mu(q)
 \left[
   \sum_{\substack{p_1,\ldots,p_7\sim Y\\
                   2p_1\cdots p_7\equiv -s\;(\bmod q)}}
   \prod_{i=1}^{7}\omega_i(p_i)
   -\mathfrak M_{7,s}(q)
 \right].
 \tag{V15-1}
\]
The required scale is $o(X/\log X)$.

\subsection*{Character parent and the separate-energy death test}

For
\[
 S_i(\chi):=\sum_{p\sim Y}\omega_i(p)\chi(p),
\]
unit orthogonality gives, after subtraction of the principal term,
\[
 D_q
 =
 \frac1{\phi(q)}
 \sum_{\substack{\chi\bmod q\\ \chi\neq\chi_0}}
 \overline{\chi(a_s)}\prod_{i=1}^{7}S_i(\chi).
 \tag{V15-2}
\]
Splitting $3+4$ and applying the ordinary multiplicative large sieve separately gives
\[
 |\mathcal B_{7,s}^{\rm off}|
 \ll
 Q^{-1}(Q^2X^{3/7})^{1/2}(Q^2X^{4/7})^{1/2}X^{o(1)}
 =
 X^{11/10+o(1)}.
 \tag{V15-3}
\]
Thus the separate-energy route misses the desired $X$ scale by $X^{1/10-o(1)}$.
The analogous separated sixth/eighth moment interpolation gives the same $X^{11/10+o(1)}$
scale, so no separated-$S_i$ high-moment argument closes this cell.

The equal-seven exponent geometry also does not satisfy the three-block range constraints of
the cited Proposition~6.3 route.  A termwise prime-box Vaughan/Heath--Brown split does not
remove this obstruction, because on a prime input at least one nonzero term retains the prime
unsplit in one factor slot.  Accordingly
\[
 \texttt{P63-BALANCED7-DICTIONARY: FALSE},
 \qquad
 \texttt{HB-PRIMEBOX-P63-UNSPLIT-NOGO45: PASS}.
 \tag{V15-4}
\]

\subsection*{Exact squarefree polarization}

For $\mathbf z=(z_1,\ldots,z_7)\in\mathbb T^7$ define, on primes $p\sim Y$,
\[
 a_{\mathbf z}(p):=\frac17\sum_{i=1}^{7}z_i\omega_i(p),
 \qquad |a_{\mathbf z}(p)|\le1,
 \tag{V15-5}
\]
and define the squarefree-supported multiplicative function
\[
 f_{\mathbf z}(p)=a_{\mathbf z}(p)\quad(p\sim Y),\qquad
 f_{\mathbf z}(p^\ell)=0\quad(\ell\ge2),
 \tag{V15-6}
\]
extended multiplicatively and set to zero if a prime factor lies outside the actual prime box.

For distinct primes $p_1,\ldots,p_7$,
\[
 [z_1\cdots z_7]\,f_{\mathbf z}(p_1\cdots p_7)
 =
 7^{-7}
 \sum_{\sigma\in S_7}
 \prod_{j=1}^{7}\omega_{\sigma(j)}(p_j).
 \tag{V15-7}
\]
Consequently the ordered labelled squarefree source is recovered exactly as the torus
coefficient
\[
 \begin{aligned}
 &\sum_{\substack{p_1,\ldots,p_7\sim Y\\
                  p_i\ {\rm distinct}\\
                  2p_1\cdots p_7\equiv-s\;(\bmod q)}}
 \prod_{i=1}^{7}\omega_i(p_i)
 \\
 &\qquad =
 7^7
 \int_{\mathbb T^7}
 \overline{z_1\cdots z_7}
 \sum_{\substack{n\asymp X\\n\equiv a_s\;(\bmod q)}}
 f_{\mathbf z}(n)W(n/X)\,d\mathbf z .
 \end{aligned}
 \tag{V15-8}
\]
This is a finite coefficient-extraction identity, not a heuristic symmetrisation.

If two labelled primes coincide, at most six prime variables remain free.  Divisor
multiplicity for $q\mid 2p_1\cdots p_7+s$ costs only $X^{o(1)}$, so
\[
 \mathcal E_{\rm rep}\ll Y^{6+o(1)}=X^{6/7+o(1)}
 =o_A\!\left(\frac{X}{(\log X)^A}\right).
 \tag{V15-9}
\]
Hence
\[
 \texttt{BALANCED7-SQUAREFREE-POLARIZATION45: PASS},
 \qquad
 \texttt{BALANCED7-REPEATED-PRIME45: PASS}.
\]

\subsection*{Smooth multiplicative class}

The Euler factor of $f_{\mathbf z}$ is $1+a_{\mathbf z}(p)p^{-w}$.  Therefore
\[
 \Lambda_{f_{\mathbf z}}(p^\ell)
 =
 (-1)^{\ell-1}a_{\mathbf z}(p)^\ell\log p,
\]
and
\[
 |\Lambda_{f_{\mathbf z}}(p^\ell)|\le \log p=\Lambda(p^\ell).
 \tag{V15-10}
\]
Thus uniformly in $\mathbf z$,
\[
 f_{\mathbf z}\in\mathcal C,
\]
in the standard smooth-supported multiplicative-function sense, and its support is
$Y$-smooth with $Y=X^{1/7}$.
We record
\[
 \texttt{BALANCED7-SMOOTH-MULTIPLICATIVE-DICTIONARY45: PASS}.
 \tag{V15-11}
\]

Published smooth-supported distribution technology lies close to this point, but the
literal endpoint dictionary needed here is not supplied by the statements used in the
programme: the $3/5-\varepsilon$ theorem requires a sufficiently small smoothness exponent
without an explicit threshold proving $1/7$ admissible, while the stronger smooth-number
range is stated with a sufficiently large parameter $C$ and does not literally certify
$C\le7$.  These are range/source mismatches, not proofs of impossibility.

\subsection*{Polarized discrepancy identity and the two remaining V-branch pins}

Let $\mathcal E_q(f_{\mathbf z};a_s)$ denote the principal plus low-conductor expected term
in the chosen beyond-$1/2$ convention, and define
\[
 \Delta_{\rm LC}(f_{\mathbf z};X,q,a_s)
 :=
 \sum_{\substack{n\asymp X\\n\equiv a_s\;(\bmod q)}}
 f_{\mathbf z}(n)W(n/X)
 -
 \mathcal E_q(f_{\mathbf z};a_s).
 \tag{V15-12}
\]
Define the coefficient-extracted expected term
\[
 \mathfrak M^{\rm pol}_{7,s}(q)
 :=
 7^7\int_{\mathbb T^7}
 \overline{z_1\cdots z_7}\,
 \mathcal E_q(f_{\mathbf z};a_s)\,d\mathbf z.
 \tag{V15-13}
\]
Then the exact source reduction is
\[
 \begin{aligned}
 \mathcal B_{7,s}(X)
 ={}&
 7^7\int_{\mathbb T^7}
 \overline{z_1\cdots z_7}
 \sum_{q\sim X^{3/5}}
 \mu(q)\,
 \Delta_{\rm LC}(f_{\mathbf z};X,q,a_s)
 \,d\mathbf z
 \\
 &+
 \sum_{q\sim Q}\mu(q)
 \bigl(\mathfrak M^{\rm pol}_{7,s}(q)-\mathfrak M_{7,s}(q)\bigr)
 +O(X^{6/7+o(1)}).
 \end{aligned}
 \tag{V15-14}
\]

Accordingly the first exact analytic residual on the balanced-seven V-branch is
\[
 \boxed{\texttt{AFFINE287-POLARIZED-OMEGA7-SIGNED-EOD45}},
\]
namely
\[
 \boxed{
 \int_{\mathbb T^7}
 \overline{z_1\cdots z_7}
 \sum_{q\sim X^{3/5}}\mu(q)
 \Delta_{\rm LC}(f_{\mathbf z};X,q,-s\overline2)
 \,d\mathbf z
 =
 o(X/\log X).
 }
 \tag{V15-15}
\]
This asks only for the degree-$(1,\ldots,1)$ torus coefficient of the signed modulus
average; it is strictly weaker than first proving a uniform absolute Bombieri--Vinogradov
bound for every $\mathbf z$.

The source-matched comparison pin is
\[
 \boxed{\texttt{AFFINE287-MULOG-COMPARISON-LOWCOND-MATCH45}},
\]
requiring
\[
 \boxed{
 \sum_{q\sim Q}\mu(q)
 \bigl(\mathfrak M^{\rm pol}_{7,s}(q)-\mathfrak M_{7,s}(q)\bigr)
 =
 o(X/\log X).
 }
 \tag{V15-16}
\]
If (V15-15) and (V15-16) hold, then by (V15-14)
\[
 \mathcal B_{7,s}(X)=o(X/\log X),
\]
so the balanced-seven modulus-average cell is closed.

\begin{status}[V15 controlling balanced-seven verdict]
\[
\begin{array}{ll}
\texttt{BALANCED7-CHARACTER-PARENT45} & \text{PASS},\\
\texttt{BALANCED7-SEPARATE-MULT-LS} & \text{NONCLOSING: }X^{1/10}\text{ deficit},\\
\texttt{P63-BALANCED7-DICTIONARY} & \text{FALSE},\\
\texttt{BALANCED7-SQUAREFREE-POLARIZATION45} & \text{PASS},\\
\texttt{BALANCED7-REPEATED-PRIME45} & \text{PASS: }X^{6/7+o(1)},\\
\texttt{BALANCED7-SMOOTH-MULTIPLICATIVE-DICTIONARY45} & \text{PASS},\\
\texttt{AFFINE287-POLARIZED-OMEGA7-SIGNED-EOD45} & \textbf{FIRST EXACT V-BRANCH ANALYTIC OPEN},\\
\texttt{AFFINE287-MULOG-COMPARISON-LOWCOND-MATCH45} & \text{SOURCE OPEN},\\
\texttt{AFFINE287-BALANCED7-MODULUS-AVERAGE45} & \text{REDUCED, NOT CLOSED},\\
\texttt{AFFINE287-PRIME-MODULUS-MU-TWOOUTER45} & \text{M/Vaughan BRANCH OPEN},\\
\texttt{FCL} & \text{OPEN},\\
\text{Erd\H{o}s \#287} & \text{OPEN}.
\end{array}
\]
\end{status}
